% Çeviri yapay zekâ tarafından üretilmiştir. Bu eşitleme 36 numaralı makalenin matematiksel metnini değiştirmez. Kaynak sunumu: OpenAI Codex — GPT-6 Astra, Ultra akıl yürütme düzeyi. Eserlerdeki yazar adları korunur.
% Haklar ilgili hak sahiplerinde kalır. Ticari olmayan, yalnızca eğitim amaçlı ücretsiz dağıtım, hak sahiplerinin itiraz etmeyeceği varsayımına dayanır; kesinleşmiş bir izin değildir. Ücretsiz erişim türev yayına izin vermez. Hak sahipleri aşağıdan itiraz edebilir; yayıncı istek üzerine ilgili malzemeyi kaldırır. Yeni lisans veya kamu malına bırakma beyanı yoktur.
% https://github.com/KokunoYumeto/modern-latex-manuscripts/issues/new
\documentclass[11pt]{article}
\usepackage[a4paper,margin=24mm]{geometry}
\usepackage{fontspec}
\usepackage{polyglossia}
\setdefaultlanguage{turkish}
\setmainfont{Noto Serif}
\setlength{\parindent}{1.2em}
\setlength{\parskip}{0.5em}
\setlength{\emergencystretch}{2em}

\begin{document}

\section*{36. İdeallerin türevlenmesi ve different}

\begin{center}
\emph{J. Ber. d. DMV 39 (1930), s. 17}
\end{center}

2. E. Noether, Göttingen: İdeallerin türevlenmesi ve different.

Bir cebirsel sayı cisminin different adı verilen idealinin, bu sayı cismini tanımlayan bir \emph{idealin} \emph{türevi} olarak nasıl ele alınabileceği gösterilmektedir. Bu tanımın alışılagelmiş tanımla aynı olduğu, kendi başlarına da ilginç olan ve benzer anlamda Lagrange interpolasyon formülünün bir genellemesini oluşturan yapı teoremlerinden çıkar.

Ayrıntılı bir sunumun \emph{Math. Ann.} dergisinde yayımlanması öngörülmektedir.

\end{document}
